\chapter{Search without a model: grid, random, quasi-random, evolutionary}
\label{sec:modelfree}

The simplest tuners evaluate configurations chosen without reference to any
model of the objective. They are worth understanding not because we would
ship them as the product's centerpiece but because they are the baselines
every fancier method must beat, they parallelize perfectly, and one of them
(random search) is quietly load-bearing inside several of the modern methods
in later chapters.

\section{Why random beats grid: an argument in one picture}

Grid search evaluates the Cartesian product of a few values per axis. Its
appeal is legibility: the results form a table a human can read. Its flaw is
arithmetic. With a budget of $n$ trials spread over $d$ axes, each axis
receives only $n^{1/d}$ distinct values. In the running example's full
fifteen-dimensional space, a generous budget of $n = 100$ trials gives each
axis $100^{1/15} \approx 1.36$ distinct values: the ``grid'' cannot even
afford two settings per knob. Even on the two-dimensional slice we can draw,
$n = 9$ trials buys a $3 \times 3$ grid, three distinct learning rates in
total.

Why is that so bad? \citet{bergstra2012random} made the now-standard
observation that machine-learning objectives typically have low
\emph{effective} dimensionality: a couple of hyperparameters dominate, but
which ones dominate varies by problem. Suppose, as is typical for PPO, the
learning rate matters enormously and the entropy coefficient barely at all.
Then what determines a trial's quality is essentially its learning-rate
coordinate alone, and the right question is not ``how well do the trials
cover the plane'' but ``how well do their \emph{projections} cover the
learning-rate axis.'' Figure~\ref{fig:gridrandom} shows the answer. The
$3 \times 3$ grid projects onto three distinct learning rates, each tested
three redundant times; nine random points project onto nine distinct
learning rates. The grid wastes two thirds of its budget re-measuring values
of the axis that matters while random search, which draws each coordinate
afresh for every trial, places $n$ distinct values on every axis
simultaneously, without needing to know in advance which axis is the
important one.

\begin{figure}[t]
\centering
\begin{tikzpicture}[scale=0.85, font=\small]
% left panel: grid
\begin{scope}
\draw[thick] (0,0) rectangle (4,4);
\foreach \x in {1,2,3}
  \foreach \y in {1,2,3}
    \fill (\x,\y) circle (2.2pt);
% projections
\foreach \x in {1,2,3}
  \fill[blue!70] (\x,-0.55) circle (2.2pt);
\draw[-{Latex}] (2,-0.05) -- (2,-0.35);
\node[anchor=north] at (2,-0.85) {3 distinct values};
\node[anchor=south] at (2,4.05) {grid, $n=9$};
\node[rotate=90, anchor=south] at (-0.35,2) {entropy coeff.\ (barely matters)};
\node[anchor=north] at (2,-1.35) {learning rate (matters)};
\end{scope}
% right panel: random
\begin{scope}[xshift=7cm]
\draw[thick] (0,0) rectangle (4,4);
\foreach \p in {(0.5,1.1),(1.1,3.4),(1.6,0.6),(2.05,2.5),(2.5,3.8),(2.9,1.7),(3.3,0.3),(3.7,2.9),(0.85,2.1)}
  \fill \p circle (2.2pt);
\foreach \x in {0.5,1.1,1.6,2.05,2.5,2.9,3.3,3.7,0.85}
  \fill[blue!70] (\x,-0.55) circle (2.2pt);
\draw[-{Latex}] (2,-0.05) -- (2,-0.35);
\node[anchor=north] at (2,-0.85) {9 distinct values};
\node[anchor=south] at (2,4.05) {random, $n=9$};
\node[anchor=north] at (2,-1.35) {learning rate (matters)};
\end{scope}
\end{tikzpicture}
\caption{The projection argument for random search, after
\citet{bergstra2012random}, drawn schematically. When one axis dominates the
objective, only the coverage of that axis's projection (blue) matters. The
same nine trials give the important axis three distinct values under a grid
and nine under random sampling, and random search does not need to know
which axis is important to collect this advantage.}
\label{fig:gridrandom}
\end{figure}

The same point can be made with elementary probability, and the number is
worth carrying around. Suppose ``good'' configurations occupy some fraction
$p$ of the search space, say the top 5 percent. Each uniform random trial
lands in the good region with probability $p$ independently, so the chance
that $n$ trials all miss it is $(1-p)^n$, and
\begin{equation}
\Pr[\text{at least one good trial}] \;=\; 1 - (1-p)^n .
\label{eq:coverage}
\end{equation}
With $p = 0.05$, sixty trials give $1 - 0.95^{60} \approx 0.95$: sixty
random trials find a top-5-percent configuration with 95 percent
probability, in any number of dimensions.

The dimension-free character of
\eqref{eq:coverage} is the honest content of the phrase ``random search
scales'': what it cannot do is exploit structure, so when the good region
is much smaller than 5 percent, or when we want the best of the good region
rather than a member of it, the model-based methods of
Chapter~\ref{sec:bayesian} earn their complexity.

The empirical side of
\citet{bergstra2012random} showed random search matching or beating grid
search on neural-network tuning at a fraction of the cost, and the result
has aged well: fifteen years later, tuning studies in reinforcement
learning still reach for random search as the honest cheap baseline
\citep{eimer2023hyperparameters}, and it remains the recommended floor in
methodological guides \citep{bischl2023hyperparameter}.

\section{Two free upgrades}

Two refinements cost nothing and should be our defaults wherever
plain random sampling appears. First, sample in a transformed space:
learning rates and other scale parameters should be drawn log-uniformly, since
the difference between $10^{-3}$ and $10^{-4}$ matters as much as the
difference between $10^{-4}$ and $10^{-5}$. A uniform draw from
$[10^{-5}, 10^{-2}]$ puts 90 percent of its samples above $10^{-3}$,
crowding into the top decade of the range; a log-uniform draw spends a
third of its budget in each of the three decades. The transform matters
more than any sophistication layered on top of it, and the same logic will
reappear as input warping inside Gaussian-process surrogates
(Chapter~\ref{sec:bayesian}).

Second, replace independent draws with a low-discrepancy sequence (Sobol
points), which spreads samples more evenly than chance and removes the
unlucky clusters and gaps of independent sampling: independent uniform draws
are perfectly happy to put two of nine samples almost on top of each other,
which purchases the same information twice. The gain is modest but free, and
quasi-random initialization is already the default first phase inside
several Bayesian optimization packages.

\section{Respecting the baseline}

Random search has one property that deserves respect when we evaluate our own
product: it is unbiased and easy to reason about. \citet{wan2022bgpbt} found
random search over their full 15-dimensional joint space to be, in their
words, a surprisingly strong baseline on several Brax environments, on par
with or better than the original population-based training on four of seven
tasks; that is our running example's own search space, and the lesson
transfers directly.

When a proposed method fails to beat well-configured
random search at matched compute, the method, not the baseline, is on trial.
That said, model-based methods do beat it when configured sensibly: the
analysis of the 2020 black-box optimization challenge, run on real tuning
problems, found Bayesian optimization entries clearly ahead of random search
across the board \citep{turner2021bayesian}.

Both facts fit in one sentence:
random search is hard to beat when its budget is generous relative to the
difficulty of the space, and clearly beatable when trials are scarce enough
that each one has to be chosen with care. Scarce, expensive trials are our
situation, which is why Chapter~\ref{sec:bayesian} exists.

\section{Adapting the sampler without modeling the objective}

Evolutionary and direct-search methods occupy a middle ground: no fitted
model of the objective, but the sampling distribution adapts. CMA-ES
(covariance matrix adaptation evolution strategy)
maintains a multivariate Gaussian over a continuous search space and, at
each generation, samples a population from it, ranks the samples by
objective value, and updates the Gaussian toward the successful ones: the
mean moves to a weighted average of the best samples, and the covariance is
updated so that directions in which good steps were recently taken become
directions the sampler explores more boldly \citep{hansen2016cmaes}.

Only
the \emph{ranking} of samples enters the update, never the raw objective
values, which buys invariance to any monotone rescaling of the objective, a
genuinely useful robustness when losses are heavy-tailed.

The price is
appetite: covariance adaptation needs generations of populations to learn
the geometry, so CMA-ES is a strong choice when trials are cheap enough to
afford hundreds to thousands of evaluations, when the space is purely
continuous, and when noise is moderate \citep{hansen2016cmaes}. In the
running example's regime of hours-long trials it is mis-matched; in the
Hamiltonian-network regime of minutes-long trials it is a live option.

Differential evolution, the other family member we will need, proposes new
points by recombining coordinate differences of existing population members:
a mutant is formed by adding the scaled difference of two population members
to a third, then crossed with a parent coordinate by coordinate. Its role
for us is as the engine inside DEHB (Chapter~\ref{sec:multifidelity}),
currently one of the most reliable budget-constrained tuners in
reinforcement learning \citep{eimer2023hyperparameters}. The practical
reading: evolutionary machinery earns its place at low fidelities and inside
hybrid methods, not as a standalone tuner for expensive trials, where its
appetite for evaluations collides with our budgets.

\section{The baseline's case, argued from four angles}
\label{sec:mf0-angles}

The chapter's verdict, transforms and Sobol sampling on by default,
random search enthroned as the mandatory baseline, grid retired to a
debugging view, the evolutionary methods wrapped or embedded rather
than built, is modest, which is why it is worth noticing that four
separate arguments underwrite it.

From arithmetic: the two computations of this chapter are one line
each and dimension-free. The exponent $n^{1/d}$ says a grid starves
every axis long before budgets we can afford; the bound
\eqref{eq:coverage} says sixty random trials find the top twentieth
with 95 percent confidence in any dimension. Between them they fix
both the floor's strength and its limit, since nothing in
\eqref{eq:coverage} improves when the good region shrinks or when we
want its best member rather than any member.

From evidence: the floor held up in the two most adversarial arenas
this document has. In the BG-PBT authors' own fifteen-dimensional
space, random search matched or beat the original population methods
on four of seven tasks (Section~\ref{sec:pop-evidence}); in the
outsiders' challenge on real tuning problems, the model-based entries
beat it across the board (Section~\ref{sec:bo-evidence}). The two
results are one fact seen from its two sides: generous budgets make
the floor hard to beat, scarce trials make deliberation pay, and our
studies live on the scarce side.

From operations: random search has no state, no hyperparameters worth
arguing about, no failure modes, and perfect parallelism, which makes
it the one component whose behavior never needs explaining; the log
and Sobol upgrades are free and are inherited by every later method,
since quasi-random points are also the GP searcher's opening phase;
and grid's single virtue, a table a human can read, survives intact
as the two-knob visual sweep.

From decision theory: the baseline's real product is epistemic. Every
verdict in the chapters ahead has the form ``this method beats the
honest cheap alternative at matched budget,'' and that sentence is
only checkable if the honest cheap alternative is permanently
installed in the validation suite. Carrying it costs nothing;
dropping it would convert every later claim from a measurement into
an opinion.

The falsification here runs in both directions. The chapter's own
positive claims are the cheapest in the document to test, the same
searcher with and without the log transform on a pinned task settles
the transform's value in an afternoon, and the chapter's main export
is the standing criterion the rest of the document is judged by: any
searcher that cannot beat the upgraded floor at matched budget, with
seeds and uncertainty reported, is not shipped.

\begin{decisions}
\noindent\textbf{Decisions from this chapter.} For each method: what it
is in one sentence, and whether it goes in the package.
\begin{itemize}[leftmargin=1.3em, itemsep=2pt]
\item \textbf{Log-scale and Sobol sampling} (include, days of work).
Drawing scale-type knobs on a log scale, and drawing the opening batch
from a scrambled Sobol sequence rather than independently. An
engineering default, not a theorem: the log transform applies to knobs
that span decades (learning rates, coefficients, step sizes) and not to
ordinals or categories, and Sobol's even coverage is a claim about
initialization batches, where it is also the GP searcher's opening
phase anyway. Two caveats worth carrying: scrambling matters, since an
unscrambled sequence's low-order points are visibly structured, and at
budgets of a handful of trials the difference from independent sampling
is below the noise we are measuring against.
\item \textbf{Random search} (include, trivial). Not as a product
feature but as the mandatory baseline in our validation suite: any
searcher that cannot beat it at matched budget is not shipped. The
arithmetic of \eqref{eq:coverage} is why it is harder to beat than its
reputation suggests.
\item \textbf{Grid search} (exclude as a search method). Keep only as a
debugging utility for two-knob visual sweeps; Figure~\ref{fig:gridrandom}
is the entire case against it as a tuner.
\item \textbf{CMA-ES} (include by wrapping an existing implementation,
low priority). Only competitive when a study can afford hundreds of
cheap, continuous trials, which for us means the Hamiltonian-network
regime and nothing else. Building it ourselves would buy nothing.
\item \textbf{Differential evolution} (no standalone entry). It enters
the package only inside DEHB in Chapter~\ref{sec:multifidelity}.
\end{itemize}
\end{decisions}

\section{What would overturn these verdicts}

The floor baseline is structural, not empirical: if the GP searcher
cannot beat log-scale Sobol at matched budget on the pinned tasks of
Chapter~\ref{sec:validation}, it is demoted and this chapter
rewritten. CMA-ES's inclusion is conditional on the wrap remaining
thin; if the dependency becomes unmaintained or the adapter requires
ongoing surgery, the verdict changes to exclude and the register
records it. Grid search's debug-view limit holds unless a colleague
requests it as a search method with a use case the decision boxes
did not consider, at which point the box is reopened.
